\documentclass[11pt]{article}
\usepackage[a4paper,margin=2.4cm]{geometry}
\usepackage{amsmath,amssymb,amsthm}
\usepackage{mathtools}
\usepackage{enumitem}
\usepackage[colorlinks=true,linkcolor=blue,urlcolor=blue]{hyperref}
\newcommand{\R}{\mathbb{R}}
\newcommand{\E}{\mathbb{E}}
\newcommand{\ii}{\mathrm{i}}
\newcommand{\dd}{\,\mathrm{d}}
\newcommand{\ch}{\varphi}
\newtheorem{remark}{Remark}

\title{Lecture 07 --- The Heston Model: CIR Variance, Affine Structure, and Simulation}
\author{Computational Finance (WI4151), Cluster 4 --- Models}
\date{}

\begin{document}
\maketitle

\noindent
These notes cover the Heston stochastic volatility model: the SDE system, the CIR
variance process and the Feller condition, the affine machinery that produces the
characteristic function, and the simulation schemes (Euler with truncation/reflection,
exact Broadie--Kaya). The single most important object for this course is the
\emph{Heston characteristic function} of $\ln S_T$, because it is exactly the input the
COS method needs to price under Heston. Cross-reference: \texttt{notes/Lecture06-notes.pdf}
(the COS pricing formula and the cumulant-based truncation range $[a,b]$).

\section{Why stochastic volatility}

Black--Scholes assumes a constant volatility $\sigma$. Empirically, option-implied
volatility depends on strike and maturity (the volatility \emph{smile}/\emph{skew}), so a
single $\sigma$ cannot reprice the whole surface. Two standard cures: add jumps to the
geometric Brownian motion (Lecture~8), or make the volatility itself stochastic. The
Heston model takes the second route. Its appeal over competitors such as SABR
($\dd S_t = \sigma_t S_t^\beta \dd W_t$, $\dd \sigma_t = \alpha \sigma_t \dd Z_t$) is that
it admits a \emph{semi-analytic} pricing solution (through its characteristic function)
and an \emph{exact} simulation scheme. SABR relies on the Hagan asymptotic formula, which
is only accurate for a limited range of maturities.

\section{The Heston SDE system}

Under the risk-neutral measure $\mathbb{Q}$ the Heston model is the two-dimensional
diffusion
\begin{equation}
\dd S_t = r S_t \dd t + \sqrt{\nu_t}\, S_t \dd W^S_t,
\qquad
\dd \nu_t = \kappa(\bar\nu - \nu_t)\dd t + \gamma \sqrt{\nu_t}\, \dd W^\nu_t,
\label{eq:heston}
\end{equation}
with the two Brownian motions \emph{correlated},
\begin{equation}
\dd W^S_t\, \dd W^\nu_t = \rho_{S,\nu}\dd t.
\end{equation}
Here $\nu_t$ is the instantaneous \emph{variance} (not the volatility; the volatility is
$\sqrt{\nu_t}$). The parameters are:
\begin{itemize}[noitemsep]
\item $\bar\nu$ --- the long-run mean of the variance;
\item $\kappa>0$ --- the speed of mean reversion of $\nu_t$ to $\bar\nu$;
\item $\gamma>0$ --- the volatility of variance (``vol-of-vol''), which governs the
curvature of the smile;
\item $\rho_{S,\nu}$ --- the spot/variance correlation, which governs the \emph{slope}
(skew) of the implied-volatility curve. As $\rho_{S,\nu}$ becomes more negative the skew
steepens; this is the leverage effect.
\end{itemize}
The variance follows a Cox--Ingersoll--Ross (CIR) process: mean-reverting with a
diffusion coefficient $\gamma\sqrt{\nu_t}$ that vanishes as $\nu_t\to0$, which is what
keeps the process non-negative.

\paragraph{Log-spot form via Cholesky.}
For pricing and simulation it is convenient to (i) move to $X_t=\ln S_t$ and (ii)
decorrelate the Brownian drivers. By It\^o's lemma on $X_t=\ln S_t$, with
$\dd S_t/S_t = r\dd t + \sqrt{\nu_t}\dd W^S_t$ and $(\dd S_t/S_t)^2=\nu_t\dd t$,
\begin{equation}
\dd \ln S_t = \Big(r - \tfrac12 \nu_t\Big)\dd t + \sqrt{\nu_t}\,\dd W^S_t.
\end{equation}
Now write the correlated pair $(W^S,W^\nu)$ in terms of two \emph{independent}
Brownian motions $(\widetilde W_1, \widetilde W_2)$. For the $2\times2$ correlation
matrix
\[
\Sigma = \begin{pmatrix} 1 & \rho \\ \rho & 1 \end{pmatrix}
= L L^\top,
\qquad
L = \begin{pmatrix} 1 & 0 \\ \rho & \sqrt{1-\rho^2}\end{pmatrix},
\]
($\rho\equiv\rho_{S,\nu}$) the Cholesky construction $W = L\widetilde W$ gives
$W^S_t = \widetilde W_1(t)$ and
$W^\nu_t = \rho\,\widetilde W_1(t) + \sqrt{1-\rho^2}\,\widetilde W_2(t)$.
One checks $\operatorname{Cov}(W^S_t,W^\nu_t)=\rho\,\E[(\widetilde W_1)^2]=\rho t$, hence
correlation $\rho$ as required. Inverting so that the \emph{variance} is driven by its
own clean increment $W^\nu$, the equivalent system is
\begin{align}
\dd \ln S_t &= \Big(r - \tfrac12\nu_t\Big)\dd t
+ \sqrt{\nu_t}\Big(\rho_{S,\nu}\dd W^S_t + \sqrt{1-\rho_{S,\nu}^2}\,\dd W^\nu_t\Big),\\
\dd \nu_t &= \kappa(\bar\nu-\nu_t)\dd t + \gamma\sqrt{\nu_t}\,\dd W^\nu_t,
\label{eq:hestonlog}
\end{align}
where now $W^S$ and $W^\nu$ on the right are \emph{independent}. This is the form used
for both Euler simulation and the affine derivation.

\section{The CIR variance process}

\subsection{Conditional distribution}

A central fact: conditional on $\nu_t$, the future variance $\nu_T$ has a known
distribution. Writing $\tau=T-t$,
\begin{equation}
\nu_T \mid \nu_t \;\sim\; \frac{Y}{2c},
\qquad
c = \frac{2\kappa}{\gamma^2\big(1-e^{-\kappa\tau}\big)},
\end{equation}
where $Y$ follows a \emph{noncentral chi-squared} distribution $\chi^{\star}_d(\lambda)$
with
\begin{equation}
d = \frac{4\kappa\bar\nu}{\gamma^2} \text{ degrees of freedom},
\qquad
\lambda = 2c\,e^{-\kappa\tau}\nu_t \text{ noncentrality}.
\end{equation}
Recall $\chi^\star_d(\lambda)$ is the law of $\sum_{i=1}^d X_i^2$ where the $X_i$ are
independent $N(\mu_i,1)$ and $\lambda=\sum_i\mu_i^2$. Its density is
\begin{equation}
f(x;d,\lambda) = \tfrac12 e^{-(x+\lambda)/2}
\Big(\tfrac{x}{\lambda}\Big)^{d/4-1/2}
I_{d/2-1}\!\big(\sqrt{\lambda x}\big),
\label{eq:ncx2}
\end{equation}
with $I_\nu$ the modified Bessel function of the first kind,
$I_\nu(y)=(y/2)^\nu\sum_{j\ge0}\frac{(y^2/4)^j}{j!\,\Gamma(\nu+j+1)}$.

\subsection{The Feller condition}

\begin{remark}[Feller condition]
If
\begin{equation}
2\kappa\bar\nu \ge \gamma^2
\end{equation}
then $\nu_t$ stays strictly positive almost surely (the origin is inaccessible). If it is
violated, $\nu_t$ can reach $0$ (and, for the pure CIR process, instantaneously reflects /
leaves $0$ again; in either case $\nu_t$ never goes negative).
\end{remark}

The cleanest way to see why this is the relevant threshold is through the density
\eqref{eq:ncx2}. As $x\to0^+$, $I_\nu(\sqrt{\lambda x})\sim (\sqrt{\lambda x}/2)^\nu/\Gamma(\nu+1)$,
so the density behaves like $x^{d/4-1/2}\cdot x^{(d/2-1)/2}=x^{d/2-1}$ up to constants
(combining the explicit power and the Bessel power). For the density not to blow up at the
origin we need the exponent $\tfrac{d}{2}-1\ge0$, i.e.
\begin{equation}
q_F \coloneqq \frac{2\kappa\bar\nu}{\gamma^2}-1 \ge 0
\iff 2\kappa\bar\nu\ge\gamma^2,
\end{equation}
recovering the Feller condition. Interpretation: a large vol-of-vol $\gamma$ relative to
the restoring force $\kappa\bar\nu$ makes excursions to zero likely; a strong restoring
force keeps the variance away from the origin. This matters in practice because the
non-central chi-squared exact scheme is only guaranteed well-defined when the Feller
condition holds (it is a \emph{sufficient} condition; rare exceptions exist but are not of
practical interest).

\section{Affine structure and the characteristic function}

\subsection{Affine diffusions: the general result}

An $n$-dimensional diffusion $\dd X_t = \mu(X_t)\dd t + \sigma(X_t)\dd W_t$ is
\emph{affine} if the drift and the diffusion covariance are affine in the state:
\begin{align}
\mu(x) &= K_0 + K_1 x, \qquad K_0\in\R^n,\ K_1\in\R^{n\times n},\\
\big(\sigma(x)\sigma(x)^\top\big)_{ij} &= (H_0)_{ij} + (H_1)_{ij}\cdot x,
\qquad H_0\in\R^{n\times n},\ H_1\in\R^{n\times n\times n},
\end{align}
and the short rate is affine, $R(x)=\rho_0+\rho_1\cdot x$. The key theorem (Duffie--Pan--Singleton 2000) states that the discounted transform
\begin{equation}
\psi^X(u,x,t,T) = \E\!\left[\exp\!\Big(-\!\int_t^T R(X_s)\dd s\Big)\,e^{u\cdot X_T}
\,\Big|\,\mathcal F_t\right]
= e^{\alpha(u,t,T)+\beta(u,t,T)\cdot x},
\label{eq:affineans}
\end{equation}
i.e. the transform is \emph{exponential-affine} in the state. The coefficients
$\alpha,\beta$ solve a system of complex ODEs (the Riccati system). For a pure diffusion
(no jumps) the system is
\begin{align}
\dot\beta &= \rho_1 - K_1^\top\beta - \tfrac12\beta^\top H_1\beta, \label{eq:riccati-b}\\
\dot\alpha &= \rho_0 - K_0\cdot\beta - \tfrac12\beta^\top H_0\beta, \label{eq:riccati-a}
\end{align}
with boundary conditions $\beta(u,T,T)=u$ and $\alpha(u,T,T)=0$. Here
$\beta^\top H_1\beta$ is the $n$-vector with $k$-th component
$\sum_{ij}\beta_i (H_1)_{ijk}\beta_j$.

\paragraph{Where the ODEs come from.}
The Feynman--Kac theorem says $\psi^X$ (as a function of $(t,x)$) satisfies the backward
PDE
\[
\partial_t \psi + \mathcal L\psi - R(x)\psi = 0,
\qquad \psi(T,x)=e^{u\cdot x},
\]
where $\mathcal L$ is the generator
$\mathcal L g = \sum_i \mu_i \partial_{x_i} g + \tfrac12\sum_{ij}(\sigma\sigma^\top)_{ij}\partial_{x_i x_j}g$.
Substituting the ansatz $\psi=e^{\alpha+\beta\cdot x}$ gives
$\partial_{x_i}\psi=\beta_i\psi$, $\partial_{x_ix_j}\psi=\beta_i\beta_j\psi$, and
$\partial_t\psi=(\dot\alpha+\dot\beta\cdot x)\psi$. Dividing by $\psi$ and collecting the
terms that are constant in $x$ versus linear in $x$ (using that $\mu$, $\sigma\sigma^\top$
and $R$ are affine) decouples into the two ODE families \eqref{eq:riccati-b}--\eqref{eq:riccati-a}: the linear-in-$x$ part gives $\dot\beta$, the constant part gives $\dot\alpha$. The
affine structure is exactly what makes the $x$-dependence separate cleanly.

\subsection{Is Heston affine? Restoring it in the log-domain}

In the $(S,\nu)$ coordinates Heston is \emph{not} affine: the diffusion covariance is
\[
\sigma\sigma^\top =
\begin{pmatrix} \nu_t S_t^2 & \gamma\rho_{S,\nu}\nu_t S_t \\
\gamma\rho_{S,\nu}\nu_t S_t & \gamma^2\nu_t \end{pmatrix},
\]
whose top-left entry $\nu_t S_t^2$ is not affine in $(S_t,\nu_t)$. The fix is to switch to
$X_t=\ln S_t$. With state $(\ln S_t,\nu_t)^\top$ the drift becomes
\[
\begin{pmatrix} \dd\ln S_t\\ \dd\nu_t\end{pmatrix}
=\Bigg(
\underbrace{\begin{pmatrix} r\\ \kappa\bar\nu\end{pmatrix}}_{K_0}
+\underbrace{\begin{pmatrix} 0 & -1/2\\ 0 & -\kappa\end{pmatrix}}_{K_1}
\begin{pmatrix}\ln S_t\\ \nu_t\end{pmatrix}
\Bigg)\dd t + \sqrt{\nu_t}\,L\,\dd\widetilde W_t,
\]
which is affine in the drift, and the covariance becomes
\[
\sigma\sigma^\top = \nu_t
\begin{pmatrix} 1 & \gamma\rho_{S,\nu}\\ \gamma\rho_{S,\nu} & \gamma^2\end{pmatrix},
\]
affine in $\nu_t$. Matching the template: $H_0=0$ (the $2\times2$ zero matrix), and $H_1$
is the $2\times2\times2$ tensor whose first slice (coefficient of $\ln S$) is the zero
matrix and whose second slice (coefficient of $\nu$) is
$\big(\begin{smallmatrix} 1 & \gamma\rho_{S,\nu}\\ \gamma\rho_{S,\nu} & \gamma^2\end{smallmatrix}\big)$.
So in log-coordinates Heston \emph{is} affine, and the general transform result applies.

\subsection{The variance ODE solved explicitly (warm-up)}

Before the full Heston cf, it is instructive to solve the one-dimensional CIR Riccati
exactly, which is what underlies the variance distribution. For the cf of $\nu_T$ alone,
$\psi^\nu(u,\nu_t,t,T)=\E[e^{u\nu_T}\mid\nu_t]=e^{\alpha_1(t)+\beta_1(t)\nu_t}$, with
$K_0=\kappa\bar\nu$, $K_1=-\kappa$, $H_0=0$, $H_1=\gamma^2$, $\rho_0=\rho_1=0$. Using the
reparametrization $\psi^\nu=e^{\alpha(\tau)+\beta(\tau)\nu_t}$ in $\tau=T-t$, the Riccati
equation reads
\begin{equation}
\dot\beta = -\kappa\beta + \tfrac12\gamma^2\beta^2,
\qquad \beta(0)=u.
\end{equation}
This is solvable by the substitution $g=1/\beta$, which linearizes it:
$\dot g = -\dot\beta/\beta^2 = \kappa g - \tfrac12\gamma^2$. The integrating factor
$e^{-\kappa t}$ gives $\dd(g e^{-\kappa t}) = -\tfrac12\gamma^2 e^{-\kappa t}\dd t$, hence
\[
g(t)e^{-\kappa t}-g(0) = -\frac{\gamma^2}{2}\int_0^t e^{-\kappa s}\dd s,
\qquad
g(t)=e^{\kappa t}\Big(g(0)-\tfrac{\gamma^2(1-e^{-\kappa t})}{2\kappa}\Big),
\]
with $g(0)=1/u$. Inverting,
\begin{equation}
\beta(t) = \frac{u\,e^{-\kappa t}}{1-\tfrac{u\gamma^2}{2}\xi_\kappa(t)},
\qquad
\xi_\kappa(t)=\frac{1-e^{-\kappa t}}{\kappa}.
\end{equation}
Then $\dot\alpha=\kappa\bar\nu\,\beta$ integrates (noting $e^{-\kappa s}=\xi_\kappa'(s)$)
to
\begin{equation}
\alpha(t)= \kappa\bar\nu u\int_0^t
\frac{e^{-\kappa s}}{1-\tfrac{u\gamma^2}{2}\xi_\kappa(s)}\dd s
= -\frac{2\kappa\bar\nu}{\gamma^2}
\ln\!\Big(1-\tfrac{u\gamma^2}{2}\xi_\kappa(t)\Big).
\end{equation}
Collecting,
\begin{equation}
\E[e^{u\nu_T}\mid\nu_t]
= \Big(1-\tfrac{u\gamma^2}{2}\xi_\kappa(\tau)\Big)^{-2\kappa\bar\nu/\gamma^2}
\exp\!\Big(\frac{u\,e^{-\kappa\tau}}{1-\tfrac{u\gamma^2}{2}\xi_\kappa(\tau)}\,\nu_t\Big).
\end{equation}
Although derived for real $u$, this extends to complex $u$ within the convergence domain,
and is exactly the cf of the (scaled) noncentral chi-squared variable above. This is the
same Riccati mechanism that produces the full Heston cf below.

\subsection{The Heston characteristic function}

Now take the full log-Heston system, set $u=\ii\omega$ purely imaginary, and write
$y_t=\ln S_t$. The transform
$\E[e^{u y_T}\mid y_t,\nu_t]=e^{\alpha(\tau,u)+u y_t+\beta(\tau,u)\nu_t}$ has, with
$\tau=T-t$ and the shorthands
\begin{equation}
b = \gamma\rho_{S,\nu}u-\kappa,
\qquad a = u(1-u),
\qquad s = \sqrt{b^2 + a\gamma^2},
\end{equation}
the closed forms
\begin{align}
\beta(\tau,u) &= \frac{a\,(1-e^{-s\tau})}{2s-(s+b)(1-e^{-s\tau})}, \label{eq:hbeta}\\
\alpha(\tau,u) &= -r\tau + ru\tau
- \kappa\bar\nu\left(\frac{s+b}{\gamma^2}\,\tau
+ \frac{2}{\gamma^2}\ln\!\Big[1-\frac{s+b}{2s}\big(1-e^{-s\tau}\big)\Big]\right).
\label{eq:halpha}
\end{align}
The Heston \emph{characteristic function} of the log-price is then
\begin{equation}
\boxed{\;
\ch_{\ln S_T}(\omega; \tau) \;=\; \E\!\big[e^{\ii\omega\ln S_T}\mid \ln S_t,\nu_t\big]
\;=\; \exp\!\big(\alpha(\tau,\ii\omega) + \ii\omega\,\ln S_t + \beta(\tau,\ii\omega)\,\nu_t\big).
\;}
\label{eq:hestoncf}
\end{equation}
This is the load-bearing object of the lecture. Everything --- semi-analytic pricing,
calibration, exact simulation of $\int\nu\,\dd s$ --- runs through it.

\begin{remark}[The ``little Heston trap'' / branch-cut subtlety]
The square root $s=\sqrt{b^2+a\gamma^2}$ and the complex logarithm in
\eqref{eq:halpha} are multi-valued. The way the formula is grouped
matters. The original Heston (1993) presentation uses an algebraically equivalent
grouping (with a factor $g=(b+s)/(b-s)$ and $\ln(1-g e^{s\tau})$ with $e^{+s\tau}$ rather
than $e^{-s\tau}$). For long maturities $\tau$ that form repeatedly crosses the branch cut
of the principal-branch $\ln$, the function becomes discontinuous and the integrand
oscillates wrongly, giving \emph{wildly wrong prices}. Albrecher et al.\ (2007) showed
that the mathematically equivalent ``little trap'' formulation --- precisely the one in
\eqref{eq:halpha}, written with $e^{-s\tau}$ and $1-\tfrac{s+b}{2s}(1-e^{-s\tau})$ ---
stays on the principal branch and is numerically stable for all $\tau$, provided one takes
the principal square root with non-negative real part. The lesson: the two forms are equal
as analytic functions but not as numerical recipes; always use the stable grouping.
\end{remark}

\subsection{The pricing PDE (consistency check)}

It\^o on $e^{-rt}V(t,S,\nu)$ and setting the drift to zero (so the discounted price is a
martingale) gives the Heston pricing PDE
\begin{equation}
\partial_t V + rS\,\partial_S V + \kappa(\bar\nu-\nu)\partial_\nu V
+ \tfrac12\nu S^2 \partial_{SS}V + \tfrac12\gamma^2\nu\,\partial_{\nu\nu}V
+ \gamma\nu S\rho_{S,\nu}\,\partial_{S\nu}V - rV = 0,
\end{equation}
where the cross term used $\dd S\,\dd\nu = \gamma\nu S\rho_{S,\nu}\dd t$. The cf
\eqref{eq:hestoncf} is the fundamental solution of the corresponding backward equation for
the transform; this is the same PDE the Feynman--Kac/affine derivation invokes.

\section{How the Heston cf plugs into COS}

Recall from Lecture~06 the COS European pricing formula. With $x=\ln(S_0/K)$, density
support truncated to $[a,b]$, and cosine coefficients $V_k$ of the payoff,
\begin{equation}
v(x,t_0) \approx e^{-r\tau}\sideset{}{'}\sum_{k=0}^{N-1}
\operatorname{Re}\!\Big\{\ch\!\Big(\tfrac{k\pi}{b-a};x\Big)\,
e^{-\ii k\pi\frac{a}{b-a}}\Big\}\,V_k,
\label{eq:cos}
\end{equation}
where ${\sum}'$ halves the first term. The \emph{only} model-specific input is the
characteristic function $\ch(\cdot)$ of the log-price. For Black--Scholes one substitutes
the Gaussian cf; for Heston one substitutes exactly \eqref{eq:hestoncf}, evaluated at the
COS frequencies $\omega_k = k\pi/(b-a)$:
\begin{equation}
\ch\!\Big(\tfrac{k\pi}{b-a};x\Big)
= \exp\!\Big(\alpha\big(\tau,\tfrac{\ii k\pi}{b-a}\big)
+ \tfrac{\ii k\pi}{b-a}\,x
+ \beta\big(\tau,\tfrac{\ii k\pi}{b-a}\big)\nu_0\Big).
\end{equation}
This is why the affine machinery matters for this course: COS is a quadrature of the
inverse Fourier transform, and Heston is exactly the model whose density is unknown in
closed form but whose cf is known in closed form, which is the regime where COS dominates.

\begin{remark}[Truncation range under Heston]
The cumulant-based rule from Lecture~06,
$[a,b]=[c_1 - L\sqrt{c_2+\sqrt{c_4}},\; c_1 + L\sqrt{c_2+\sqrt{c_4}}]$, requires the
cumulants $c_1,c_2,c_4$ of $\ln S_T$, which for Heston follow by differentiating
$\log\ch$ at $\omega=0$. The danger I hit in Assignment~2 is generic: at \emph{short}
maturities $\tau$ the variance of $\ln S_T$ is small (small $c_2$), so $[a,b]$ becomes
narrow; if the higher cumulants or a fat left tail (negative $\rho$) are under-weighted,
the range can chop off mass and COS loses its spectral accuracy. Under Heston the skew
from $\rho<0$ and the vol-of-vol $\gamma$ inflate the tails, so the prudent fix is a
larger $L$ (e.g.\ $L=10$--$12$) and including $c_4$, rather than trusting the BS-style
$L=8$ to $10$ with $c_4$ ignored.
\end{remark}

\section{Monte Carlo simulation under Heston}

\subsection{Euler scheme and the negative-variance problem}

The naive Euler discretization of \eqref{eq:hestonlog} on a grid $t_i$ is
\begin{align}
\nu_{t_{i+1}} &= \nu_{t_i} + \kappa(\bar\nu-\nu_{t_i})(t_{i+1}-t_i)
+ \gamma\sqrt{\nu_{t_i}}\sqrt{t_{i+1}-t_i}\,Z,\\
\ln S_{t_{i+1}} &= \ln S_{t_i} + \Big(r-\tfrac12\nu_{t_i}\Big)(t_{i+1}-t_i)
+ \sqrt{\nu_{t_i}}\sqrt{t_{i+1}-t_i}\,Y,
\end{align}
with $Y=\rho_{S,\nu}Z + \sqrt{1-\rho_{S,\nu}^2}\,X$, $X,Z\sim N(0,1)$ independent.
Because $Z$ is Gaussian and unbounded below, $\mathbb P(\nu_{t_{i+1}}<0)>0$ always: the
discretized variance can go negative, after which $\sqrt{\nu}$ is undefined. No naive
time-stepping scheme avoids this.

\paragraph{Quick fixes (with errors).}
\begin{itemize}[noitemsep]
\item \emph{Truncation / floor}: set $\nu_{t_{i+1}}\leftarrow\max(\nu_{t_{i+1}},\varepsilon)$
for tiny $\varepsilon\ge0$. Simple, but introduces a truncation bias.
\item \emph{Reflection}: set $\nu_{t_{i+1}}\leftarrow|\nu_{t_{i+1}}|$. Useful when the
Feller condition fails; the reflected path systematically sits a little above the floored
one near zero.
\item \emph{Log-variance change of variable}: simulate $Y_t=\ln\nu_t$ by Euler and
exponentiate. This does \emph{not} work: $\ln\nu$ can blow up, so $\nu=e^{Y}$ explodes for
some paths.
\end{itemize}
These are all biased; the bias is worst exactly when the Feller condition is badly
violated and the variance spends time near $0$.

\subsection{Exact simulation of the variance}

A bias-free alternative for the variance uses the conditional law directly:
$\nu_{t_{i+1}}\mid\nu_{t_i}\sim Y/(2c)$ with $Y\sim\chi^\star_d(\lambda)$,
$c=2\kappa/[\gamma^2(1-e^{-\kappa\Delta})]$, $d=4\kappa\bar\nu/\gamma^2$,
$\lambda=2c\,e^{-\kappa\Delta}\nu_{t_i}$, $\Delta=t_{i+1}-t_i$. Sampling the noncentral
chi-squared uses:
\begin{itemize}[noitemsep]
\item if $d>1$: $\chi^\star_d(\lambda)^2 \stackrel{d}{=} (Z+\sqrt\lambda)^2 + \chi^2_{d-1}$
with $Z\sim N(0,1)$ and an independent ordinary $\chi^2_{d-1}$;
\item for general $d>0$: draw $J\sim\mathrm{Poisson}(\lambda/2)$, then conditionally
$\chi^\star_d(\lambda)^2\sim\chi^2_{d+2J}$.
\end{itemize}
This is exact for $\nu$ and produces strictly non-negative samples. Caveat: the noncentral
chi-squared representation needs $d>0$, i.e.\ effectively the Feller condition; when it
fails one switches to the log-variance transform of Fang--Oosterlee (2011).

\subsection{Exact (Broadie--Kaya) simulation of the spot}

Knowing the variance law is not enough for $S$, because the log-price depends on
$\int_u^t\nu_s\dd s$ and on the stochastic integrals against $W^\nu,W^S$. Integrating
\eqref{eq:heston} over $[u,t]$,
\begin{align}
S_t &= S_u \exp\!\Big(r(t-u) - \tfrac12\!\int_u^t\!\nu_s\dd s
+ \rho_{S,\nu}\!\int_u^t\!\sqrt{\nu_s}\dd W^\nu_s
+ \sqrt{1-\rho_{S,\nu}^2}\!\int_u^t\!\sqrt{\nu_s}\dd W^S_s\Big),\\
\nu_t &= \nu_u + \kappa\bar\nu(t-u) - \kappa\!\int_u^t\!\nu_s\dd s
+ \gamma\!\int_u^t\!\sqrt{\nu_s}\dd W^\nu_s.
\end{align}
The second line lets us express one stochastic integral in terms of the others. The
Broadie--Kaya (2006) algorithm is then:
\begin{enumerate}[label=\textbf{Step \arabic*.},leftmargin=*]
\item Sample $\nu_t\mid\nu_u$ from the noncentral chi-squared above.
\item Sample $\int_u^t\nu_s\dd s$ given $\nu_u,\nu_t$ --- \emph{the hard step}.
\item Recover $\int_u^t\sqrt{\nu_s}\dd W^\nu_s
= \tfrac1\gamma\big(\nu_t-\nu_u-\kappa\bar\nu(t-u)+\kappa\int_u^t\nu_s\dd s\big)$.
\item Sample $\int_u^t\sqrt{\nu_s}\dd W^S_s \mid \int_u^t\nu_s\dd s
\sim N\!\big(0,\int_u^t\nu_s\dd s\big)$ (independence of $W^S$ from the variance path).
\end{enumerate}
Then assemble $S_t$. The result is unbiased and lets you jump directly from $u$ to $t$ with
no intermediate time steps.

\paragraph{Step 2 via the cf of the integrated variance.}
The distribution of $\int_u^t\nu_s\dd s$ given $(\nu_u,\nu_t)$ is obtained from its
characteristic function $\ch(\omega)=\E^{\mathbb Q}[\exp(\ii\omega\int_u^t\nu_s\dd s)\mid\nu_u,\nu_t]$,
which is known in closed form (a ratio of modified Bessel functions $I_{d/2-1}$, with
$R(\omega)=\sqrt{\kappa^2-2\gamma^2\ii\omega}$ and $d=4\kappa\bar\nu/\gamma^2$). One then
inverts to get the CDF $F$ --- the original paper used the trapezoidal rule on the Fourier
inversion; \emph{today one uses the COS method} --- and finally samples by inversion
$\int\nu\,\dd s = F^{-1}(U)$ with $U\sim\mathrm{Uniform}(0,1)$, valid because
$\mathbb P(F^{-1}(U)\le r)=\mathbb P(U\le F(r))=F(r)$.

\paragraph{Trade-offs.} Exact simulation has no discretization/truncation error, higher
accuracy per path, and direct sampling at horizon $t$; but it is slow and intricate to
implement (Bessel functions, Fourier inversion per step). Broadie--Kaya's tables show the
exact scheme's RMS error decaying as $O(N^{-1/2})$ in the number of trials with zero bias,
versus Euler's slower convergence dominated by discretization bias.

\subsection{The three tools behind the integrated-variance cf (sketch)}

The closed form for $\ch(\omega)$ is built from three ingredients; knowing the names and
roles is enough for the oral.
\begin{itemize}[noitemsep]
\item \textbf{Infinitesimal generator.} For $\dd X_t=\mu(X_t)\dd t+\sigma(X_t)\dd W_t$ the
generator is $\mathcal A f = \mu f' + \tfrac12\sigma^2 f''$. The CIR generator is
$\mathcal A f(x)=\kappa(\bar\nu-x)f'(x)+\tfrac12\gamma^2 x f''(x)$. The $n$-dimensional
\emph{squared Bessel} process $\dd X_t = n\dd t + 2\sqrt{X_t}\dd W_t$ has generator
$nf'(x)+2xf''(x)$ and a known Laplace transform of its time integral (Pitman--Yor 1982).
\item \textbf{Time change.} The scaling $cW_t\sim W_{c^2 t}$ lets us rescale the diffusion
coefficient. The time change $s=\tfrac{\gamma^2}{4}t$ maps the CIR second-order coefficient
$\tfrac12\gamma^2 x$ onto the squared-Bessel value $2x$.
\item \textbf{Change of measure (Girsanov).} After the time change the \emph{drift}
coefficients still differ ($\kappa(\bar\nu-x)$ vs $n$); a Girsanov change of measure absorbs
the discrepancy, expressing the CIR integrated-variance Laplace transform as a ratio of
squared-Bessel Laplace transforms, which Pitman--Yor give explicitly. Matching coefficients
yields the Bessel-ratio formula for $\ch(\omega)$.
\end{itemize}

\section{Parameter intuition (for the smile)}

\begin{itemize}[noitemsep]
\item $\gamma$ (vol-of-vol): with $\rho=0$, larger $\gamma$ deepens the smile curvature.
\item $\rho_{S,\nu}$: more negative $\rho$ steepens the skew (left tail heavier).
\item $\kappa$: limited effect on the smile shape ($1$--$2\%$); controls how fast ATM
implied vol converges to $\sqrt{\bar\nu}$ in maturity.
\item $\nu_0,\bar\nu$: similar effects, mediated by $\kappa$ ($\nu_0$ dominates short
maturities, $\bar\nu$ long maturities).
\end{itemize}
Versus Black--Scholes, Heston gives heavier tails of log returns and reproduces the smile.

\section*{Honest gaps / things not on these slides}

\begin{itemize}[noitemsep]
\item The \emph{full} explicit Riccati derivation of \eqref{eq:hbeta}--\eqref{eq:halpha}
(as opposed to the scalar CIR warm-up) is not carried out term by term on the slides; they
quote the closed form. I have shown the mechanism on the 1-D CIR case, which is the same.
\item The little-Heston-trap discussion (branch cut, Albrecher et al.) is \emph{my
addition}; the slides give the stable grouping \eqref{eq:halpha} but do not flag the trap
explicitly. Worth knowing for Fang's exam.
\item The exact integrated-variance cf (the Bessel-ratio formula, slide 53) and the full
Girsanov derivation (slides 65--69) are stated; I have summarized them rather than
reproduced every line. The squared-Bessel Laplace transform itself (Pitman--Yor) is quoted
without proof.
\item Cumulant formulas $c_1,c_2,c_4$ for Heston needed for the COS truncation range are
referenced from Lecture~06 / Fang--Oosterlee; not derived here.
\end{itemize}

\end{document}
